\hspace{20pt}
\begin{center}
\textbf{{\huge{Аннотация}}}
\end{center}
\vspace{8pt}

Выпускная квалификационная работа посвящена исследованию методов идентификации объектов при съёмке с разноракурсных камер наблюдения на примере задачи реидентификации транспортных средств (Vehicle ReID). В рамках работы формализуется постановка задачи ReID как поиска соответствий между запросом (query) и множеством кандидатов (gallery) при условии разнесённости идентичностей между обучением и тестом. Рассматриваются особенности «городского» видеонаблюдения: сильные изменения ракурса и масштаба, вариации освещения, частичные окклюзии, низкое разрешение и доменный сдвиг между камерами. В исследовании исключается использование государственных регистрационных знаков; идентификация выполняется по внешнему виду.

В качестве эмпирической базы применяются общепринятые наборы данных для Vehicle ReID (например, VeRi-776, VehicleID; дополнительно для стресс‑тестов — VERI‑Wild / CityFlowV2). Предложен и реализован воспроизводимый пайплайн: детекция/кроп объекта, извлечение признаков нейросетевой моделью, сопоставление в эмбеддинг‑пространстве и постобработка результатов. В качестве базовых моделей реализованы и сравниваются две архитектуры — CNN (ResNet-50) и Transformer (ViT-Small) — с BNNeck и функциями потерь Cross‑Entropy (с ослаблением меток) и Triplet Loss; исследуются эффекты стратегий обучения (PK‑семплинг, аугментации Random Erasing/Color Jitter) и постобработки (k‑reciprocal re‑ranking).

Качество методов оценивается метриками ранжирования CMC@k и среднеарифметической точностью (mAP). Практическая значимость работы состоит в построении инженерно‑простого и переносимого решения ReID для транспортных средств в городских условиях без опоры на распознавание номеров. Результаты и код организованы так, чтобы обеспечить воспроизводимость и дальнейшее расширение.

\noindent\textbf{Ключевые слова:} реидентификация, Vehicle ReID, многокамерное наблюдение, эмбеддинги, Triplet Loss, BNNeck, re‑ranking, mAP, CMC.
